
\clearpage
\section{Training Data Details}\label{appendix:data}

\begin{table}[th]
    \small
    \centering
    \adjustbox{width=\textwidth,center}{
    \begin{tabular}{m{1cm}<{\centering}cr|m{1cm}<{\centering}cr|m{1cm}<{\centering}cr}
    \toprule
        \tb{ISO Code} & \tb{Language} & \tb{Samples} & \tb{ISO Code} & \tb{Language} & \tb{Samples} & \tb{ISO Code} & \tb{Language} & \tb{Samples} \\
    \cmidrule(r){1-3}\cmidrule(lr){4-6}\cmidrule(l){7-9}
eng & English & 16,059,324 & msa & Malay & 175,111 & tat & Tatar & 22,327 \\
zho & Chinese & 4,280,372 & hrv & Croatian & 146,132 & bos & Bosnian & 21,175 \\
rus & Russian & 3,426,943 & isl & Icelandic & 124,495 & yor & Yoruba & 20,139 \\
spa & Spanish & 2,771,907 & slv & Slovenian & 122,065 & min & Minangkabau & 19,868 \\
fra & French & 2,426,075 & srp & Serbian & 113,663 & che & Chechen & 19,518 \\
deu & German & 1,749,922 & urd & Urdu & 113,258 & arz & Egyptian Arabic & 16,783 \\
ara & Arabic & 1,402,943 & ben & Bengali & 87,787 & lmo & Lombard & 16,575 \\
nld & Dutch & 1,395,135 & aze & Azerbaijani & 82,209 & arg & Aragonese & 16,500 \\
vie & Vietnamese & 1,159,472 & afr & Afrikaans & 81,233 & bak & Bashkir & 16,451 \\
hin & Hindi & 1,106,611 & tam & Tamil & 78,384 & som & Somali & 16,369 \\
kor & Korean & 1,083,205 & kat & Georgian & 77,567 & als & Tosk Albanian & 15,655 \\
jpn & Japanese & 1,082,466 & tel & Telugu & 77,362 & ido & Ido & 15,613 \\
ita & Italian & 960,595 & mal & Malayalam & 76,518 & szl & Silesian & 14,845 \\
ind & Indonesian & 952,218 & mon & Mongolian & 59,851 & wuu & Wu Chinese & 14,762 \\
por & Portuguese & 919,257 & nno & Norwegian Nynorsk & 58,298 & new & Nepal Bhasa & 14,714 \\
pol & Polish & 885,709 & kaz & Kazakh & 55,317 & chv & Chuvash & 13,759 \\
tur & Turkish & 694,051 & cym & Welsh & 53,951 & pnb & Western Panjabi & 13,657 \\
dan & Danish & 675,313 & mar & Marathi & 53,803 & fry & Western Frisian & 13,649 \\
tha & Thai & 654,534 & sqi & Albanian & 53,602 & snd & Sindhi & 13,210 \\
swe & Swedish & 577,848 & nob & Norwegian Bokmål & 52,905 & ori & Oriya & 12,791 \\
fas & Persian & 535,107 & pus & Pushto & 52,753 & plt & Plateau Malagasy & 12,717 \\
ukr & Ukrainian & 471,079 & mkd & Macedonian & 52,504 & scn & Sicilian & 12,694 \\
ces & Czech & 467,569 & hbs & Serbo-Croatian & 48,523 & kur & Kurdish & 11,872 \\
nor & Norwegian & 424,995 & ceb & Cebuano & 47,408 & sun & Sundanese & 11,793 \\
ell & Modern Greek & 378,202 & jav & Javanese & 47,283 & bar & Bavarian & 11,193 \\
cat & Catalan & 370,156 & war & Waray (Philippines) & 45,348 & yid & Yiddish & 10,785 \\
ron & Romanian & 344,845 & kan & Kannada & 44,534 & ckb & Central Kurdish & 9,829 \\
fin & Finnish & 332,394 & epo & Esperanto & 44,266 & fao & Faroese & 9,825 \\
bul & Bulgarian & 321,379 & lat & Latin & 43,335 & ina & Interlingua & 9,782 \\
tgl & Tagalog & 313,985 & guj & Gujarati & 40,184 & gla & Scottish Gaelic & 9,769 \\
glg & Galician & 301,386 & uzb & Uzbek & 39,951 & bug & Buginese & 9,662 \\
mya & Burmese & 294,167 & amh & Amharic & 38,763 & que & Quechua & 9,406 \\
hye & Armenian & 288,622 & oci & Occitan & 37,413 & bpy & Bishnupriya & 9,400 \\
khm & Khmer & 287,530 & bel & Belarusian & 33,330 & san & Sanskrit & 8,730 \\
nep & Nepali & 276,057 & azb & South Azerbaijani & 31,815 & lim & Limburgan & 8,573 \\
hun & Hungarian & 271,802 & kir & Kirghiz & 29,319 & hau & Hausa & 8,435 \\
eus & Basque & 270,551 & mlg & Malagasy & 28,661 & mai & Maithili & 8,180 \\
heb & Hebrew & 263,869 & vol & Volapük & 27,187 & zsm & Standard Malay & 8,179 \\
lao & Lao & 244,750 & ast & Asturian & 26,004 & ibo & Igbo & 8,132 \\
swa & Swahili & 241,497 & pan & Panjabi & 25,096 & vec & Venetian & 8,121 \\
azj & North Azerbaijani & 213,046 & ltz & Luxembourgish & 25,092 & ilo & Iloko & 7,968 \\
lav & Latvian & 212,058 & nds & Low German & 24,713 & asm & Assamese & 7,042 \\
sin & Sinhala & 207,903 & hat & Haitian & 23,940 & sah & Yakut & 7,011 \\
slk & Slovak & 202,049 & bre & Breton & 23,931 & arb & Standard Arabic & 6,945 \\
tgk & Tajik & 200,631 & gle & Irish & 23,148 & sna & Shona & 6,933 \\
est & Estonian & 191,063 & sco & Scots & 23,032 & mlt & Maltese & 6,911 \\
lit & Lithuanian & 184,391 & xho & Xhosa & 22,799 & zul & Zulu & 6,669 \\
    \bottomrule
    \end{tabular}
    }
    \caption{Natural language distribution in the training data of F2LLM-v2 (part1).}
    \label{tab:language-distribution-1}
\end{table}

\begin{table}[th]
    \small
    \centering
    \adjustbox{width=\textwidth,center}{
    \begin{tabular}{m{1cm}<{\centering}cr|m{1cm}<{\centering}cr|m{1cm}<{\centering}cr}
    \toprule
        \tb{ISO Code} & \tb{Language} & \tb{Samples} & \tb{ISO Code} & \tb{Language} & \tb{Samples} & \tb{ISO Code} & \tb{Language} & \tb{Samples} \\
    \cmidrule(r){1-3}\cmidrule(lr){4-6}\cmidrule(l){7-9}
mzn & Mazanderani & 6,352 & tsn & Tswana & 1,539 & lvs & Standard Latvian & 800 \\
uig & Uighur & 6,190 & mwl & Mirandese & 1,491 & mag & Magahi & 800 \\
oss & Iron Ossetic & 5,893 & div & Dhivehi & 1,387 & mni & Manipuri & 800 \\
tuk & Turkmen & 5,854 & kbp & Kabiyè & 1,349 & mos & Mossi & 800 \\
ary & Moroccan Arabic & 5,703 & chm & Mari & 1,238 & nqo & N'Ko & 800 \\
wln & Walloon & 5,408 & ewe & Ewe & 1,220 & nus & Nuer & 800 \\
cdo & Min Dong Chinese & 5,175 & smo & Samoan & 1,175 & ory & Odia & 800 \\
npi & Nepali & 5,156 & tso & Tsonga & 1,174 & prs & Dari & 800 \\
nap & Neapolitan & 4,778 & fij & Fijian & 1,122 & quy & Ayacucho Quechua & 800 \\
ace & Achinese & 4,758 & bam & Bambara & 1,061 & sat & Santali & 800 \\
mrj & Western Mari & 4,728 & lin & Lingala & 1,046 & tpi & Tok Pisin & 800 \\
xmf & Mingrelian & 4,712 & nav & Navajo & 1,028 & tum & Tumbuka & 800 \\
pes & Iranian Persian & 4,414 & roh & Romansh & 999 & tzm & Central Atlas Tamazight & 800 \\
diq & Dimli & 4,140 & ssw & Swati & 982 & umb & Umbundu & 800 \\
apc & Levantine Arabic & 4,079 & awa & Awadhi & 954 & uzn & Northern Uzbek & 800 \\
wol & Wolof & 4,068 & pag & Pangasinan & 954 & ydd & Eastern Yiddish & 800 \\
pbt & Southern Pashto & 3,796 & cor & Cornish & 938 & yue & Yue Chinese & 800 \\
nso & Pedi & 3,676 & dzo & Dzongkha & 928 & dyu & Dyula & 799 \\
srd & Sardinian & 3,614 & udm & Udmurt & 890 & lua & Luba-Lulua & 799 \\
ban & Balinese & 3,581 & fon & Fon & 883 & twi & Twi & 798 \\
lij & Ligurian & 3,487 & kon & Kongo & 873 & aeb & Tunisian Arabic & 793 \\
hsb & Upper Sorbian & 3,440 & glv & Manx & 864 & kik & Kikuyu & 761 \\
acq & Ta'izzi-Adeni Arabic & 3,188 & tir & Tigrinya & 851 & tyv & Tuvinian & 626 \\
crh & Crimean Tatar & 2,985 & pms & Piemontese & 842 & ava & Avaric & 588 \\
mri & Maori & 2,922 & myv & Erzya & 840 & aym & Aymara & 587 \\
egl & Emilian & 2,859 & sag & Sango & 827 & krc & Karachay-Balkar & 587 \\
ars & Najdi Arabic & 2,712 & run & Rundi & 825 & ful & Fulah & 581 \\
grn & Guarani & 2,705 & acm & Mesopotamian Arabic & 800 & orm & Oromo & 548 \\
nya & Chichewa & 2,607 & aka & Akan & 800 & stq & Saterfriesisch & 461 \\
hif & Fiji Hindi & 2,566 & ayr & Central Aymara & 800 & lah & Lahnda & 450 \\
kas & Kashmiri & 2,363 & bem & Bemba & 800 & ton & Tonga & 391 \\
fur & Friulian & 2,284 & bho & Bhojpuri & 800 & mdf & Moksha & 314 \\
swh & Swahili & 2,253 & cjk & Chokwe & 800 & haw & Hawaiian & 299 \\
fil & Filipino & 2,156 & dik & Southwestern Dinka & 800 & nia & Nias & 297 \\
sme & Northern Sami & 2,156 & fuv & Nigerian Fulfulde & 800 & bis & Bislama & 272 \\
shn & Shan & 2,098 & gaz & West Central Oromo & 800 & alt & Southern Altai & 250 \\
sot & Southern Sotho & 2,089 & hne & Chhattisgarhi & 800 & srn & Sranan Tongo & 204 \\
kin & Kinyarwanda & 2,031 & kab & Kabyle & 800 & ven & Venda & 194 \\
lug & Ganda & 1,994 & kac & Kachin & 800 & kbd & Kabardian & 172 \\
pap & Papiamento & 1,974 & kam & Kamba (Kenya) & 800 & xal & Kalmyk & 122 \\
cos & Corsican & 1,949 & kea & Kabuverdianu & 800 & din & Dinka & 104 \\
mhr & Eastern Mari & 1,633 & khk & Halh Mongolian & 800 & jam & Jamaican Creole English & 100 \\
bjn & Banjar & 1,600 & kmb & Kimbundu & 800 & kal & Kalaallisut & 92 \\
knc & Central Kanuri & 1,600 & kmr & Northern Kurdish & 800 & iku & Inuktitut & 84 \\
taq & Tamasheq & 1,600 & ltg & Latgalian & 800 & guc & Wayuu & 52 \\
kom & Komi & 1,583 & luo & Luo & 800 & chr & Cherokee & 51 \\
bod & Tibetan & 1,563 & lus & Lushai & 800 & ady & Adyghe & 33 \\
    \bottomrule
    \end{tabular}
    }
    \caption{Natural language distribution in the training data of F2LLM-v2 (part2).}
    \label{tab:language-distribution-2}
\end{table}

\begin{table}[th]
    \small
    \centering
    \begin{tabular}{cr|cr}
    \toprule
Language & Samples & Language & Samples \\
    \cmidrule(r){1-2}\cmidrule(lr){3-4}
python & 1,972,390 & css & 1,003 \\
php & 553,651 & typescript & 888 \\
java & 483,469 & r & 636 \\
cpp & 393,514 & lisp & 467 \\
go & 351,586 & jsx & 436 \\
javascript & 245,632 & objective-c & 327 \\
c\# & 92,008 & json & 264 \\
ruby & 68,317 & xml & 258 \\
c & 43,487 & yaml & 180 \\
rust & 12,924 & assembly & 171 \\
kotlin & 11,284 & powershell & 162 \\
sql & 6,826 & vba & 157 \\
pascal & 6,299 & lua & 126 \\
d & 5,278 & matlab & 114 \\
haskell & 4,967 & dart & 107 \\
scala & 4,120 & bash & 105 \\
html & 2,777 & http & 99 \\
shell & 2,095 & graphql & 89 \\
perl & 2,009 & svg & 82 \\
swift & 1,907 & vb.net & 75 \\
ocaml & 1,894 & groovy & 63 \\
csharp & 1,891 & Misc. & 2,301 \\
    \bottomrule
    \end{tabular}
    \caption{Programming language distribution in the training data of F2LLM-v2.}
    \label{tab:language-distribution-3}
\end{table}
\begin{table}[th]
    \tiny
    \centering
    \adjustbox{width=\textwidth,center}{
    \begin{tabular}{lccrm{9cm}}
    \toprule
        \tb{Name} & \tb{Language} & \tb{Format} & \tb{Size} & \tb{URL} \\
    \midrule
\rowcolor{gray!15} \multicolumn{5}{c}{Bitext Mining}\\
UNPC~\citep{2016unpc} & 6 & Retrieval & 2,922,245 & \url{huggingface.co/datasets/Helsinki-NLP/un_pc} \\
ParaCrawl~\citep{2020paracrawl} & 30 & Retrieval & 10,684,184 & \url{paracrawl.eu/index.php} \\
BactrianX Translation~\citep{2023bactrianx} & 52 & Clustering & 491,282 & \url{huggingface.co/datasets/MBZUAI/Bactrian-X} \\
Europarl~\citep{2005europarl} & 21 & Clustering & 477,566 & \url{huggingface.co/datasets/Helsinki-NLP/europarl} \\
\midrule
\rowcolor{gray!15} \multicolumn{5}{c}{Question Answering}\\
WebFAQ~\citep{2025webfaq} & 49 & Retrieval & 4,368,504 & \url{huggingface.co/datasets/PaDaS-Lab/webfaq-retrieval} \\
mMARCO~\citep{2021mmarco} & 14 & Retrieval & 5,470,174 & \url{huggingface.co/datasets/unicamp-dl/mmarco} \\
PAQ~\citep{2021PAQ} & en & Retrieval & 938,771 & \url{huggingface.co/datasets/sentence-transformers/paq} \\
SQuAD~\citep{2016SQuAD} & en & Retrieval & 89,509 & \url{huggingface.co/datasets/rajpurkar/squad} \\
Stack Exchange~\citep{2021StackExchangeDataset} & en & Retrieval & 754,705 & \url{huggingface.co/datasets/flax-sentence-embeddings/stackexchange_titlebody_best_voted_answer_jsonl} \\
Arguana~\citep{2018Arguana} & en & Retrieval & 22,848 & \url{huggingface.co/datasets/BeIR/arguana-generated-queries} \\
Natural Questions~\citep{2019NaturalQuestions} & en & Retrieval & 97,209 & \url{huggingface.co/datasets/sentence-transformers/natural-questions} \\
HotpotQA~\citep{2018HotpotQA} & en & Retrieval & 120,528 & \url{huggingface.co/datasets/mteb/hotpotqa} \\
ELI5~\citep{2019ELI5} & en & Retrieval & 161,345 & \url{huggingface.co/datasets/Pavithree/eli5} \\
FiQA2018~\citep{2018FIQA} & en & Retrieval & 7,452 & \url{huggingface.co/datasets/mteb/fiqa} \\
BioASQ~\citep{2015BioASQ} & en & Retrieval & 125,248 & \url{huggingface.co/datasets/BeIR/bioasq-generated-queries} \\
NFCorpus~\citep{2016NFCorpus} & en & Retrieval & 1,283 & \url{huggingface.co/datasets/mteb/nfcorpus} \\
TriviaQA~\citep{2017TriviaQA} & en & Retrieval & 60,025 & \url{huggingface.co/datasets/sentence-transformers/trivia-qa-triplet} \\
PubMedQA~\citep{2019PubMedQA} & en & Retrieval & 60,227 & \url{huggingface.co/datasets/qiaojin/PubMedQA} \\
Amazon QA~\citep{2019AmazonQA} & en & Retrieval & 59,340 & \url{github.com/amazonqa/amazonqa} \\
MIRACL~\citep{2023MIRACL} & 16 & Retrieval & 26,740 & \url{huggingface.co/datasets/miracl/miracl} \\
Mr.TyDi~\citep{2021mrtidy} & 11 & Retrieval & 48,619 & \url{huggingface.co/datasets/mteb/mrtidy} \\
MLDR~\citep{2024BGE-M3} & 13 & Retrieval & 40,264 & \url{huggingface.co/datasets/Shitao/MLDR} \\
MKQA~\citep{2021mkqa} & 26 & Retrieval & 69,287 & \url{huggingface.co/datasets/mteb/MKQARetrieval} \\
StackOverflowQA~\citep{2025CoIR} & en & Retrieval & 13,820 & \url{huggingface.co/datasets/mteb/stackoverflow-qa} \\
ProCQA~\citep{2025procqa} & 11 & Retrieval & 485,780 & \url{github.com/jordane95/procqa} \\
Yahoo\_Answers~\citep{2015yahooanswers} & en & Retrieval & 196,645 & \url{huggingface.co/datasets/sentence-transformers/yahoo-answers} \\
GooAQ~\citep{2021gooaq} & en & Retrieval & 473,876 & \url{github.com/allenai/gooaq} \\
T2Ranking~\citep{2023t2ranking} & zh & Retrieval & 85,521 & \url{huggingface.co/datasets/sentence-transformers/t2ranking} \\
DuReader~\citep{2018dureader} & zh & Retrieval & 78,023 & \url{huggingface.co/datasets/sentence-transformers/dureader} \\
cMedQAv2~\citep{2018cmedqav2} & zh & Retrieval & 23,105 & \url{huggingface.co/datasets/sentence-transformers/cmedqa-v2} \\
Huatuo\_kgqa~\citep{2025huatuoqa} & zh & Retrieval & 53,835 & \url{huggingface.co/datasets/FreedomIntelligence/huatuo_knowledge_graph_qa} \\
Huatuo\_encqa~\citep{2025huatuoqa} & zh & Retrieval & 253,523 & \url{huggingface.co/datasets/FreedomIntelligence/huatuo_encyclopedia_qa} \\
Multi CPR Medical~\citep{2022multicprmedical} & zh & Retrieval & 62,085 & \url{github.com/Alibaba-NLP/Multi-CPR} \\
HealthCareMagic~\citep{2023healthcaremagic} & en & Retrieval & 78,626 & \url{github.com/Kent0n-Li/ChatDoctor} \\
MedicalQA\_ru~\citep{2022medicalqaru} & ru & Retrieval & 71,932 & \url{huggingface.co/datasets/blinoff/medical_qa_ru_data} \\
LLM Retrieval Data~\citep{2024llm_retrieval_data} & zh & Retrieval & 177,850 & \url{huggingface.co/datasets/infgrad/retrieval_data_llm} \\
RefGPT~\citep{2023refgpt} & zh & Retrieval & 184,332 & \url{github.com/sufengniu/RefGPT} \\
Lawzhidao~\citep{2020lawzhidao} & zh & Retrieval & 11,899 & \url{heywhale.com/mw/dataset/5e953ca8e7ec38002d02fca7/content} \\
MedMCQA~\citep{2022MedMCQA} & en & Retrieval & 16,526 & \url{huggingface.co/datasets/openlifescienceai/medmcqa} \\
CMCQA~\citep{2022CMCQA} & zh & Retrieval & 108,529 & \url{github.com/WENGSYX/CMCQA} \\
MedQA~\citep{2020MedQA} & en, zh & Retrieval & 13,458 & \url{github.com/jind11/MedQA} \\
webMedQA~\citep{2019webMedQA} & zh & Retrieval & 27,122 & \url{github.com/hejunqing/webMedQA} \\
MedQuAD~\citep{2019MedQuAD} & en & Retrieval & 14,268 & \url{github.com/abachaa/MedQuAD} \\
Medical Flashcards~\citep{2023medical_flashcards} & en & Retrieval & 33,183 & \url{huggingface.co/datasets/medalpaca/medical_meadow_medical_flashcards} \\
MailruQA~\citep{2022mailruqa} & ru & Retrieval & 150,777 & \url{huggingface.co/datasets/Den4ikAI/mailruQA-big} \\
WikiOmnia~\citep{2022wikiomnia} & ru & Retrieval & 462,693 & \url{huggingface.co/datasets/RussianNLP/wikiomnia} \\
PersianQA~\citep{2021PersianQA} & fa & Retrieval & 6,277 & \url{github.com/SajjjadAyobi/PersianQA} \\
PQuAD~\citep{2023PQuAD} & fa & Retrieval & 46,699 & \url{github.com/AUT-NLP/PQuAD} \\
ParSQuAD~\citep{2021parsquad} & fa & Retrieval & 41,879 & \url{github.com/BigData-IsfahanUni/ParSQuAD} \\
MQA~\citep{2021mqa} & 38 & Retrieval & 5,131,895 & \url{huggingface.co/datasets/clips/mqa} \\
HotpotQA-NL~\citep{2025beir-nl} & nl & Retrieval & 81,192 & \url{huggingface.co/datasets/clips/beir-nl-hotpotqa} \\
\midrule
\rowcolor{gray!15} \multicolumn{5}{c}{Instruction Data}\\
Aya~\citep{2024aya} & 65 & Retrieval & 126,965 & \url{huggingface.co/datasets/CohereLabs/aya_dataset} \\
MURI~\citep{2025muri} & 194 & Retrieval & 720,782 & \url{huggingface.co/datasets/akoksal/muri-it} \\
OASST2~\citep{2023oasst2} & 26 & Retrieval & 12,449 & \url{huggingface.co/datasets/OpenAssistant/oasst2} \\
MultiAlpaca~\citep{2023multialpaca} & 11 & Retrieval & 125,447 & \url{huggingface.co/datasets/DAMO-NLP-MT/multialpaca} \\
WildChat~\cite{2024wildchat} & 76 & Retrieval & 638,781 & \url{huggingface.co/datasets/allenai/WildChat-4.8M} \\
M2Lingual~\citep{2025m2lingual} & 75 & Retrieval & 158,251 & \url{huggingface.co/datasets/ServiceNow-AI/M2Lingual} \\
Natural Reasoning~\citep{2025naturalreasoning} & en & Retrieval & 845,682 & \url{huggingface.co/datasets/facebook/natural_reasoning} \\
Infinity Instruct~\citep{2025infinityinstruct} & en, zh & Retrieval & 757,439 & \url{huggingface.co/datasets/BAAI/Infinity-Instruct} \\
COIG~\citep{2025coig} & zh & Retrieval & 42,415 & \url{huggingface.co/datasets/m-a-p/COIG-CQIA} \\
Medinstruct~\citep{2023medinstruct} & en & Retrieval & 51,539 & \url{github.com/XZhang97666/AlpaCare} \\
CodeFeedbackST~\citep{2025CoIR} & 137 & Retrieval & 115,971 & \url{huggingface.co/datasets/mteb/codefeedback-st} \\
CodeFeedbackMT~\citep{2025CoIR} & python & Retrieval & 52,221 & \url{huggingface.co/datasets/mteb/codefeedback-mt} \\
OpenOrca~\citep{2023openorca} & en & Retrieval & 896,450 & \url{huggingface.co/datasets/Open-Orca/OpenOrca} \\
MEDI2~\citep{2025medi2} & en & Retrieval & 668,036 & \url{huggingface.co/datasets/GritLM/MEDI2} \\
MedicalInstruction~\citep{2023MedicalInstruction} & en & Retrieval & 75,268 & \url{huggingface.co/datasets/Mohammed-Altaf/medical-instruction-120k} \\
SiberianDataset~\citep{2023SiberianDatasetXL} & ru & Retrieval & 255,663 & \url{huggingface.co/datasets/SiberiaSoft/SiberianDatasetXL} \\
Ru Instruct~\citep{2024ru-instruct-translated} & ru & Retrieval & 452,574 & \url{huggingface.co/datasets/d0rj/ru-instruct} \\
KoAlpaca-RealQA~\citep{2024koalpaca_realqa} & ko & Retrieval & 17,599 & \url{huggingface.co/datasets/beomi/KoAlpaca-RealQA} \\
KoAlpaca~\citep{2023koalpaca} & ko & Retrieval & 21,126 & \url{huggingface.co/datasets/beomi/KoAlpaca-v1.1a} \\
KoMagpie~\citep{2024komagpie} & ko & Retrieval & 428,780 & \url{huggingface.co/datasets/channelcorp/KoMagpie-raw} \\
Nordic Text Matching~\citep{2025nordic_text_matching} & da, sv, no & Retrieval & 182,485 & \url{huggingface.co/datasets/kardosdrur/synthetic-nordic-text_matching} \\
Nordic Retrieval~\citep{2025nordic_retrieval} & da, sv, no & Retrieval & 172,437 & \url{huggingface.co/datasets/kardosdrur/synthetic-nordic-retrieval} \\
Nordic Classification~\citep{2025nordic_classification} & da, sv, no & Classification & 199,280 & \url{huggingface.co/datasets/kardosdrur/synthetic-nordic-classification} \\
\midrule
\rowcolor{gray!15} \multicolumn{5}{c}{Title Matching}\\
S2ORC-Title-Abstract~\citep{2020S2ORC} & en & Retrieval & 250,000 & \url{huggingface.co/datasets/sentence-transformers/s2orc} \\
CORD 19~\citep{2020cord19} & en & Retrieval & 373,674 & \url{huggingface.co/datasets/medalpaca/medical_meadow_cord19} \\
Multi CPR ECom~\citep{2022multicprmedical} & zh & Retrieval & 90,850 & \url{github.com/Alibaba-NLP/Multi-CPR} \\
ESCI~\citep{2022esci} & en, ja, es & Retrieval & 80,468 & \url{huggingface.co/datasets/tasksource/esci} \\
CLIRMatrix~\citep{2020clirmatrix} & 137 & Retrieval & 3,275,561 & \url{github.com/ssun32/CLIRMatrix} \\
DBPedia~\citep{2017dbpedia} & en & Retrieval & 288,736 & \url{huggingface.co/datasets/BeIR/dbpedia-entity-generated-queries} \\
\midrule
\rowcolor{gray!15} \multicolumn{5}{c}{NLI}\\
SNLI~\citep{2015SNLI} & en & Retrieval & 54,585 & \url{huggingface.co/datasets/stanfordnlp/snli} \\
MNLI~\citep{2018MNLI} & en & Retrieval & 112,075 & \url{huggingface.co/datasets/nyu-mll/multi_nli} \\
ANLI~\citep{2020ANLI} & en & Retrieval & 18,801 & \url{huggingface.co/datasets/facebook/anli} \\
XNLI~\citep{2018XNLI} & 14 & Retrieval & 1,400,600 & \url{huggingface.co/datasets/mteb/xnli} \\
OCNLI~\citep{2020OCNLI} & zh & Retrieval & 6,616 & \url{huggingface.co/datasets/dirtycomputer/OCNLI} \\
CMNLI~\citep{2024cmnli} & zh & Retrieval & 113,914 & \url{huggingface.co/datasets/fenffef/cmnli} \\
MSciNLI~\citep{2024MSciNLI} & en & Retrieval & 19,185 & \url{huggingface.co/datasets/sadat2307/MSciNLI} \\
    \bottomrule
    \end{tabular}
    }
    \caption{Number of samples in our collected training dataset (part 1).}
    \label{tab:data-1}
\end{table}


\begin{table}[th]
    \tiny
    \centering
    \adjustbox{width=\textwidth,center}{
    \begin{tabular}{lccrm{9cm}}
    \toprule
        \tb{Name} & \tb{Language} & \tb{Format} & \tb{Size} & \tb{URL} \\
    \midrule
\rowcolor{gray!15} \multicolumn{5}{c}{Topic Classification}\\
Arxiv Clustering P2P~\citep{2022mteb-arxiv} & en & Clustering & 83,476 & \url{huggingface.co/datasets/mteb/raw_arxiv} \\
Arxiv Clustering S2S~\citep{2022mteb-arxiv} & en & Clustering & 83,486 & \url{huggingface.co/datasets/mteb/raw_arxiv} \\
Biorxiv Clustering P2P~\citep{2022mteb-biorxiv} & en & Clustering & 57,296 & \url{huggingface.co/datasets/mteb/raw_biorxiv} \\
Biorxiv Clustering S2S~\citep{2022mteb-biorxiv} & en & Clustering & 57,296 & \url{huggingface.co/datasets/mteb/raw_biorxiv} \\
Medrxiv Clustering P2P~\citep{2022mteb-medrxiv} & en & Clustering & 18,659 & \url{huggingface.co/datasets/mteb/raw_medrxiv} \\
Medrxiv Clustering S2S~\citep{2022mteb-medrxiv} & en & Clustering & 18,659 & \url{huggingface.co/datasets/mteb/raw_medrxiv} \\
MLSUM Clustering~\citep{2020mlsum} & de, es, fr, ru & Clustering & 325,739 & \url{huggingface.co/datasets/mteb/mlsum} \\
TwentyNewsgroups~\citep{1995TwentyNewsGroups} & en & Clustering & 11,060 & \url{huggingface.co/datasets/SetFit/20_newsgroups} \\
SIB200ClusteringS2S~\citep{2024SIB200} & 205 & Clustering & 163,302 & \url{huggingface.co/datasets/mteb/sib200} \\
Reddit Clustering P2P~\citep{2021Reddit-Clustering-P2P} & en & Clustering & 80,000 & \url{huggingface.co/datasets/sentence-transformers/reddit-title-body} \\
Reddit Clustering S2S~\citep{2021Reddit-StackExchange-Clustering-S2S} & en & Clustering & 58,141 & \url{github.com/UKPLab/TWEAC-qa-agent-selection/tree/master/data/reddit/train} \\
Stack Exchange Clustering P2P~\citep{2021StackExchange-Clustering-P2P} & en & Clustering & 80,000 & \url{huggingface.co/datasets/flax-sentence-embeddings/stackexchange_title_body_jsonl} \\
Stack Exchange Clustering S2S~\citep{2021Reddit-StackExchange-Clustering-S2S} & en & Clustering & 56,731 & \url{github.com/UKPLab/TWEAC-qa-agent-selection/tree/master/data/stackexchange/train} \\
THUCNews~\citep{2016THUCNews} & zh & Clustering & 100,000 & \url{huggingface.co/datasets/SirlyDreamer/THUCNews} \\
TNews~\citep{2020CLUE} & zh & Clustering & 49,726 & \url{huggingface.co/datasets/C-MTEB/TNews-classification} \\
CSL~\citep{2022CSL} & zh & Clustering & 100,000 & \url{huggingface.co/datasets/neuclir/csl} \\
\midrule
\rowcolor{gray!15} \multicolumn{5}{c}{Summarization}\\
XSum~\citep{2018XSum} & en & Retrieval & 184,383 & \url{huggingface.co/datasets/EdinburghNLP/xsum} \\
CNN\_DM~\citep{2015CNN-DM} & en & Retrieval & 100,000 & \url{huggingface.co/datasets/abisee/cnn_dailymail} \\
MLSUM Retreival~\citep{2020mlsum} & 5 & Retrieval & 801,159 & \url{huggingface.co/datasets/mteb/mlsum} \\
Sentence Compression~\citep{2013Sentence-Compression} & en & Retrieval & 175,477 & \url{huggingface.co/datasets/sentence-transformers/sentence-compression} \\
\midrule
\rowcolor{gray!15} \multicolumn{5}{c}{Text-to-Code}\\
OCGI~\citep{2024OCGI} & python & Retrieval & 1,052,849 & \url{huggingface.co/datasets/nvidia/OpenCodeGeneticInstruct} \\
OpenCodeReasoning-2~\citep{2025OpenCodeReasoning2} & python, cpp & Retrieval & 16,632 & \url{huggingface.co/datasets/nvidia/OpenCodeReasoning-2} \\
xCodeEval NL2Code~\citep{2024xcodeeval} & 17 & Retrieval & 51,072 & \url{huggingface.co/datasets/NTU-NLP-sg/xCodeEval} \\
CosQA~\citep{2021CosQA} & python & Retrieval & 9,409 & \url{huggingface.co/datasets/mteb/cosqa} \\
SyntheticText2SQL~\citep{2024synthetic-text-to-sql} & sql & Retrieval & 99,617 & \url{huggingface.co/datasets/mteb/synthetic-text2sql} \\
\midrule
\rowcolor{gray!15} \multicolumn{5}{c}{Code-to-Code}\\
xCodeEval Code2Code~\citep{2024xcodeeval} & 17 & Retrieval & 37,056 & \url{huggingface.co/datasets/NTU-NLP-sg/xCodeEval} \\
xCodeEval Translation~\citep{2024xcodeeval} & 11 & Clustering & 500,000 & \url{huggingface.co/datasets/NTU-NLP-sg/xCodeEval} \\
CodeSearchNet-ccr~\citep{2025CoIR} & 6 & Retrieval & 905,195 & \url{huggingface.co/datasets/CoIR-Retrieval/CodeSearchNet-ccr} \\
\midrule
\rowcolor{gray!15} \multicolumn{5}{c}{Code-to-Text}\\
CodeSearchNet~\citep{2019CodeSearchNet} & 6 & Retrieval & 936,813 & \url{huggingface.co/datasets/CoIR-Retrieval/CodeSearchNet} \\
\midrule
\rowcolor{gray!15} \multicolumn{5}{c}{Paraphrase Detection}\\
StackExchangeDupQuestions-S2S~\citep{2021Embedding-Training-Data} & en & Retrieval & 183,559 & \url{huggingface.co/datasets/sentence-transformers/stackexchange-duplicates} \\
StackExchangeDupQuestions-P2P~\citep{2021Embedding-Training-Data} & en & Retrieval & 203,060 & \url{huggingface.co/datasets/sentence-transformers/stackexchange-duplicates} \\
QQP~\citep{2019QQP} & en & Retrieval & 243,598 & \url{gluebenchmark.com/tasks} \\
StackOverflowDupQuestions~\citep{2018StackOverflowDupQuestions} & en & Retrieval & 19,847 & \url{huggingface.co/datasets/mteb/stackoverflowdupquestions-reranking} \\
PawsX~\citep{2019pawsx} & 7 & Retrieval & 216,219 & \url{huggingface.co/datasets/google-research-datasets/paws-x} \\
LCQMC~\citep{2018LCQMC} & zh & Retrieval & 167,213 & \url{huggingface.co/datasets/C-MTEB/LCQMC} \\
\midrule
\rowcolor{gray!15} \multicolumn{5}{c}{Sentiment Analysis}\\
Amazon Polarity~\citep{2013Amazon-Reviews} & en & Classification & 100,000 & \url{huggingface.co/datasets/mteb/amazon_polarity} \\
IMDb~\citep{2011IMDB} & en & Classification & 24,904 & \url{huggingface.co/datasets/mteb/imdb} \\
Toxic Conversations~\citep{2019Toxic-Conversations} & en & Classification & 49,900 & \url{huggingface.co/datasets/mteb/toxic_conversations_50k} \\
Amazon Counterfactual~\citep{2021Amazon-Counterfactual} & en, de, ja & Classification & 14,870 & \url{huggingface.co/datasets/mteb/amazon_counterfactual} \\
Waimai~\citep{2023CMTEB} & zh & Classification & 7,999 & \url{huggingface.co/datasets/C-MTEB/waimai-classification} \\
Amazon Reviews~\citep{2013Amazon-Reviews} & 6 & Clustering & 600,000 & \url{huggingface.co/datasets/mteb/amazon_reviews_multi} \\
Emotion~\citep{2018Emotion} & en & Clustering & 17,944 & \url{huggingface.co/datasets/mteb/emotion} \\
Tweet Sentiment Extraction~\citep{2020tweet-sentiment-extraction} & en & Clustering & 26,732 & \url{huggingface.co/datasets/mteb/tweet_sentiment_extraction} \\
RuSentiment~\citep{2021ru_sentiment} & ru & Clustering & 100,000 & \url{huggingface.co/datasets/MonoHime/ru_sentiment_dataset} \\
CEDR~\citep{2020cedr} & ru & Clustering & 4,376 & \url{huggingface.co/datasets/sagteam/cedr_v1} \\
\midrule
\rowcolor{gray!15} \multicolumn{5}{c}{Intent Classification}\\
Massive Intent~\citep{2023MASSIVE} & 51 & Clustering & 661,923 & \url{huggingface.co/datasets/mteb/amazon_massive_intent} \\
MTOP Intent~\citep{2021MTOP} & 6 & Clustering & 83,922 & \url{huggingface.co/datasets/mteb/mtop_intent} \\
Banking77~\citep{2020Banking77} & en & Clustering & 9,993 & \url{huggingface.co/datasets/mteb/banking77} \\
\midrule
\rowcolor{gray!15} \multicolumn{5}{c}{Domain Classification}\\
Massive Scenario~\citep{2023MASSIVE} & 51 & Clustering & 661,923 & \url{huggingface.co/datasets/mteb/amazon_massive_scenario} \\
MTOP Domain~\citep{2021MTOP} & 6 & Clustering & 83,922 & \url{huggingface.co/datasets/mteb/mtop_domain} \\
\midrule
\rowcolor{gray!15} \multicolumn{5}{c}{Language Classification}\\
BactrianX Language Classification~\citep{2023bactrianx} & 52 & Clustering & 491,405 & \url{huggingface.co/datasets/MBZUAI/Bactrian-X} \\
\midrule
\rowcolor{gray!15} \multicolumn{5}{c}{Citation Prediction}\\
S2ORC-TItle-Citation~\citep{2020S2ORC} & en & Retrieval & 132,879 & \url{huggingface.co/datasets/sentence-transformers/s2orc} \\
S2ORC-Abstract-Citation~\citep{2020S2ORC} & en & Retrieval & 231,587 & \url{huggingface.co/datasets/sentence-transformers/s2orc} \\
SPECTER~\citep{2020SPECTER} & en & Retrieval & 24,717 & \url{huggingface.co/datasets/sentence-transformers/specter} \\
\midrule
\rowcolor{gray!15} \multicolumn{5}{c}{Linguistic Acceptability}\\
MELA~\citep{2024MELA} & 10 & Classification & 40,267 & \url{huggingface.co/datasets/Geralt-Targaryen/MELA} \\
ScaLA~\citep{2023ScaLA} & 9 & Classification & 128,471 & \url{huggingface.co/datasets/alexandrainst/scala} \\
DaLA~\citep{2025DaLA} & da & Classification & 6,508 & \url{huggingface.co/datasets/giannor/dala_large} \\
\midrule
\rowcolor{gray!15} \multicolumn{5}{c}{Claim Verification}\\
FEVER~\citep{2018FEVER} & en & Retrieval & 106,605 & \url{huggingface.co/datasets/mteb/fever} \\
SciFact~\citep{2020SciFact} & en & Retrieval & 859 & \url{huggingface.co/datasets/mteb/scifact} \\
COLIEE~\citep{2022COLIEE} & en & Retrieval & 454 & \url{www.modelscope.cn/datasets/sentence-transformers/coliee} \\
FEVER-NL~\citep{2025beir-nl} & nl & Retrieval & 94,987 & \url{huggingface.co/datasets/clips/beir-nl-fever} \\
\midrule
\rowcolor{gray!15} \multicolumn{5}{c}{STS}\\
STS12~\citep{2012STS12} & en & Retrieval & 1,858 & \url{huggingface.co/datasets/mteb/sts12-sts} \\
STS22~\citep{2022STS22} & en & Retrieval & 389 & \url{huggingface.co/datasets/mteb/sts22-crosslingual-sts} \\
STSBenchmark~\citep{2021STSBenchmark} & en & Retrieval & 3,297 & \url{huggingface.co/datasets/mteb/stsbenchmark-sts} \\
STS22-Crosslingual~\citep{2022STS22} & 7 & Retrieval & 1,469 & \url{huggingface.co/datasets/mteb/sts22-crosslingual-sts} \\
BQ~\citep{2023CMTEB} & zh & Retrieval & 2,436 & \url{huggingface.co/datasets/C-MTEB/BQ} \\
QBQTC~\citep{2021qbqtc} & zh & Retrieval & 37,139 & \url{github.com/CLUEbenchmark/QBQTC} \\
SimCLUE~\citep{2022simclue} & zh & Retrieval & 213,301 & \url{github.com/CLUEbenchmark/SimCLUE} \\
LCSTS~\citep{2015LCSTS} & zh & Retrieval & 278,146 & \url{huggingface.co/datasets/hugcyp/LCSTS} \\
Nordic STS~\citep{2025nordic_sts} & da, sv, no & Retrieval & 70,617 & \url{huggingface.co/datasets/kardosdrur/synthetic-nordic-sts} \\
    \bottomrule
    \end{tabular}
    }
    \caption{Number of samples in our collected training dataset (part 2).}
    \label{tab:data-2}
\end{table}

\clearpage
\section{Details on MTEB Evaluation}\label{appendix:mteb}



The Massive Text Embedding Benchmark (MTEB) is widely recognized as the de facto standard for the comprehensive evaluation of text embedding models. Originally introduced by \citet{2023MTEB}, it was vastly expanded into the Massive Multilingual Text Embedding Benchmark (MMTEB) through a large-scale, open-science collaboration~\citep{2025MMTEB}. This community-driven effort has established a rigorous and diverse evaluation framework, encompassing over 500 quality-controlled tasks that span more than 250 languages and a wide array of domains.

The significance of MTEB lies in its unprecedented scale and diversity, which addresses the critical limitations of previous benchmarks that were often constrained to a few languages (mostly English), specific domains (e.g., news), or a single task type (e.g., retrieval). To provide a holistic assessment of a model's capabilities, MTEB organizes its evaluation tasks into ten distinct categories:
\begin{itemize}
    \item Retrieval: Assesses a model's ability to find relevant documents from a large corpus for a given query.
    \item Reranking: Measures the ability to reorder a given list of candidate documents by their relevance to a query.
    \item Classification: Evaluates performance on standard text classification tasks (e.g., sentiment analysis, topic classification).
    \item Clustering: Tests how well embeddings group semantically similar documents together.
    \item Pair Classification: Involves predicting the relationship between a pair of texts (e.g., paraphrase detection, natural language inference).
    \item Semantic Textual Similarity (STS): Measures the ability to predict the degree of semantic similarity between two sentences on a continuous scale.
    \item Bitext Mining: Assesses the ability to identify translated sentence pairs from a collection of sentences in two languages.
    \item Summarization: Evaluates the semantic similarity between a model-generated summary and a reference summary.
    \item Instruction Reranking: A more challenging reranking variant where the model must follow a detailed natural language instruction to determine relevance.
    \item Multilabel Classification: A classification variant where each document can be assigned multiple labels.
\end{itemize}

The hundreds of tasks are further organized into benchmarks, which are curated subsets of tasks grouped by language, domain, or a combination of both. This includes language-specific benchmarks such as English, Chinese, and Russian; domain-specific benchmarks such as Code and Medical; and aggregated benchmarks like Multilingual, European, and Scandinavian, which test performance across a broad and diverse set of languages. This hierarchical structure allows for both a fine-grained analysis of a model's performance on a specific language or domain and a high-level view of its overall multilingual and multi-domain capabilities.

In this work, we leverage the breadth of MTEB to provide a robust and thorough evaluation of our models. We evaluate on \tb{17 benchmarks}, totaling \tb{430 unique tasks}: Multilingual, Code, Medical, English, Russian, French, German, Polish, Dutch, Indic, Persian, Chinese, Japanese, Korean, Vietnamese, European, and Scandinavian. This extensive evaluation allows for a robust and fine-grained assessment of our models' capabilities, directly supporting our claims of multilingual inclusivity and broad domain competence.